% Bridgeland and Toda stability for the chamberwise positive half of K3 x E.
% Primary sources: Bridgeland 2007 and 2008; Toda 2008, 2009, 2012;
% Kontsevich and Soibelman 2008, 2010; Joyce and Song 2012;
% Oberdieck and Pandharipande 2016; Gritsenko and Nikulin 1998;
% Borcherds 1998.

\providecommand{\Db}{D^{b}}
\providecommand{\Coh}{\mathrm{Coh}}
\providecommand{\Stab}{\mathrm{Stab}}
\providecommand{\Mum}{\mathcal{M}}
\providecommand{\HH}{\mathbb{H}}
\providecommand{\CoHA}{\mathrm{CoHA}}
\providecommand{\BPS}{\mathrm{BPS}}
\providecommand{\DT}{\mathrm{DT}}
\providecommand{\PT}{\mathrm{PT}}
\providecommand{\kBKM}{\kappa_{\mathrm{BKM}}}
\providecommand{\kch}{\kappa_{\mathrm{ch}}}
\providecommand{\kcat}{\kappa_{\mathrm{cat}}}
\providecommand{\kfib}{\kappa_{\mathrm{fiber}}}
\providecommand{\Yp}{Y^{+}}
\providecommand{\CY}{\mathrm{CY}}
\providecommand{\ChirAlg}{\mathrm{ChirAlg}}
\providecommand{\Phifun}{\Phi}
\providecommand{\cO}{\mathcal{O}}

\section*{Bridgeland and Toda stability for $\Db(K3\times E)$}

\subsection*{Stability data and the positive sector}

Let $X=K3\times E$ and let $\Db(X)=\Db(\Coh(X))$.  A Bridgeland
stability condition on a triangulated category $\mathcal D$ is a pair
$\sigma=(Z,\mathcal P)$ consisting of a group homomorphism
$Z\colon K_0(\mathcal D)\to\mathbb C$ and a slicing
$\mathcal P(\phi)\subset\mathcal D$ satisfying the
Harder-Narasimhan axiom and the support property.  Bridgeland 2007,
Theorem 1.2, gives a local homeomorphism from each connected component
of $\Stab(\mathcal D)$ to a linear subspace of
$\mathrm{Hom}_{\mathbb Z}(K_0(\mathcal D),\mathbb C)$.

The half-plane selector is defined only relative to a stability datum.
Here $\sigma$ denotes either a genuine Bridgeland stability condition on
a constructed component or a Toda-type DT stability datum whose central
charge determines an upper half-plane.  Set
\[
  \Gamma_\sigma^{+}
  =
  \{\gamma\in K_0(\Db(X))\mid \operatorname{Im}Z_\sigma(\gamma)>0\}.
\]
The chamberwise positive half is built from semistable critical
cohomology in this sector:
\[
  \Yp(X;\sigma)
  :=
  \bigoplus_{\gamma\in\Gamma_\sigma^{+}}
  H^{*}\!\left(\Mum_\gamma^{\sigma\text{-ss}},\phi_W\right),
\]
where $\Mum_\gamma^{\sigma\text{-ss}}$ is the moduli stack of
$\sigma$-semistable objects of class $\gamma$ and $\phi_W$ is the
vanishing-cycle sheaf on the local critical chart.  For compact
$K3\times E$, the reduced obstruction theory enters the virtual trace
and the scalar partition function; it is separate from the K3 surface
stability datum.  When the orientation datum, local critical charts,
compact-support or Borel--Moore convention, and reduced obstruction
theory required by the chosen chamber are fixed, extension stacks
\[
  \Mum_{\gamma_1}\times\Mum_{\gamma_2}
  \leftarrow
  \Mum_{\gamma_1,\gamma_2}^{\mathrm{ext}}
  \rightarrow
  \Mum_{\gamma_1+\gamma_2}
\]
give the Kontsevich-Soibelman cohomological Hall convolution product.
Associativity is the Hall theorem in this oriented critical setting.
The output is an $E_1$ Hall algebra object.  Since
$\dim_{\mathbb C}X=3$, any specialised $\Phifun_3$ output produced from
this datum is $E_1$-chiral.  Braided $E_2$ structures occur only on
centre, double, or module-category layers after the relevant pairing,
completion, and action data have been constructed.

\subsection*{The K3 factor and the product slice}

For a complex algebraic K3 surface $S$, Bridgeland 2008, Theorem 1.1,
constructs a distinguished component
$\Stab^\dagger(S)\subset\Stab(\Db(S))$ covering the period domain
$\mathcal P_0^+(S)$ of the Mukai lattice
\[
  \widetilde H(S,\mathbb Z)
  \cong U^{\oplus 4}\oplus E_8(-1)^{\oplus 2}.
\]
The central charge has the form
\[
  Z_\Omega(F)=\langle \Omega,v(F)\rangle,\qquad
  v(F)=\mathrm{ch}(F)\sqrt{\mathrm{td}(S)},
\]
with $\Omega\in\widetilde H(S,\mathbb C)$ satisfying the usual period
conditions.  The stability angle $\phi(F)=\arg Z(F)/\pi$ orders the
Harder-Narasimhan factors.

On the elliptic factor, slope stability has central charge
\[
  Z_E(F)=-\deg(F)+i\,\mathrm{rk}(F)
\]
on coherent sheaves and extends to $\Db(E)$ through the standard heart.
The external product of a K3 heart and an elliptic heart gives a
decomposable large-volume test slice on $X$.  It is a product-stable
locus, separate from a constructed full Bridgeland component of
\(\Db(X)\).  Even at
the level of charges it is proper.  For an algebraic K3 surface,
$K_{\mathrm{num}}(S)\cong \widetilde H_{\mathrm{alg}}(S,\mathbb Z)$ has
rank $\rho(S)+2$, while the full Mukai lattice
$\widetilde H(S,\mathbb Z)$ has rank $24$.  The elliptic curve has
$\operatorname{rk}K_{\mathrm{num}}(E)=2$.  Hence the algebraic product
charge lattice has rank $2(\rho(S)+2)$, and the full topological product
charge lattice has rank $24\cdot 2=48$.  A decomposable product central
charge is determined by the two factor charges, whereas a general charge
in $\mathrm{Hom}(K_{\mathrm{num}}(X),\mathbb C)$ has mixed K\"unneth
components.

Toda's limit and polynomial stability conditions supply the usable
threefold stability input: perverse-coherent hearts, tilted hearts, and
polynomial central charges control DT-type moduli on Calabi-Yau
threefolds.  The full Bridgeland support-property existence theorem for
$\Db(K3\times E)$ remains a proof obligation.  The chamberwise
construction above uses the Toda and DT stability lane when the full
support-property component is unavailable.  A Hall, BKM, or
chiral-algebra input requires the additional Hall and factorisation data
stated below.

\subsection*{Wall crossing transports ordered products}

Kontsevich and Soibelman attach to a stability structure on a
3-Calabi-Yau category, with support and orientation data, an ordered
product in a completed quantum torus:
\[
  \mathcal A_\sigma
  =
  \prod_{\arg Z_\sigma(\gamma)\ \mathrm{in\ clockwise\ order}}
  \mathcal T_\gamma^{\Omega_\sigma(\gamma)}.
\]
Across a wall, the order of the rays changes and the invariants
$\Omega_\sigma(\gamma)$ change so that the completed product remains
the same element.  The equality lives in the completed torus or
completed motivic Hall algebra.  It transports ordered products; it does
canonically identify the semistable stacks in the two chambers only
after the chosen wall-crossing correspondence is supplied.
Individual cohomology groups $H^*(\Mum_\gamma^{\sigma\text{-ss}},\phi_W)$,
CoHA isomorphisms, and BKM brackets require their own comparison maps.

The reduced $K3\times E$ scalar is compact threefold enumerative data.
It is independent of the K3 surface stability datum.  With the monic
Borcherds product
\[
  D_5=64^{-1}\Delta_5,
\]
Oberdieck and Pandharipande give
\[
  \Phi_{10}^{\mathrm{OP}}
  =D_5^2
  =4096^{-1}\Delta_5^2,
  \qquad
  Z_{\DT,\mathrm{red}}(K3\times E)
  =-(\Phi_{10}^{\mathrm{OP}})^{-1}
  =-4096\,\Delta_5^{-2}.
\]
This scalar compares with the graded character in the reduced chamber.
It is a one-dimensional trace of reduced enumerative theory.  The Hall
product, Hall pairing, Drinfeld double, and BKM root bracket require
their own construction data; a scalar partition function gives a
character constraint, while an algebra presentation requires Hall
recognition data.

\subsection*{Separated invariants on $K3\times E$}

The compact threefold, Heisenberg specialisation, BKM denominator, and
K3 fibre carry separate invariants.  K\"unneth gives
$H^{0,\bullet}(K3\times E)=H^{0,\bullet}(K3)\otimes H^{0,\bullet}(E)$,
hence
\[
  \kcat(K3\times E)
  =
  \chi(\cO_{K3\times E})
  =
  \chi(\cO_{K3})\,\chi(\cO_E)
  =
  2\cdot 0
  =
  0,
\]
\[
  \kch(K3\times E)
  =
  \sum_q(-1)^q h^{0,q}(K3\times E)
  =
  1-1+1-1
  =
  0,
\]
\[
  \kch^{\mathrm{Heis}}(K3\times E)=3,\qquad
  \kfib(K3\times E)=24,
\]
and
\[
  \kBKM(\Phi_N)=\frac{c_N(0)}2,\qquad
  \bigl(c_N(0)\bigr)_{N=1,2,3,4,6}=(10,8,6,4,2),
\]
so
\[
  \bigl(\kBKM(\Phi_N)\bigr)_{N=1,2,3,4,6}=(5,4,3,2,1).
\]
At $N=1$, $\Delta_5$ has weight $5$ and $c_1(0)=10$.  The Igusa square
$\Phi_{10}^{\mathrm{OP}}$ has weight $10$ in the OP scalar
normalisation, while the primitive BKM denominator has weight
$\kBKM(\Delta_5)=5$.

These are the primitive CHL rows.  The $N\in\{5,7,8\}$ diagonal-divisor
rows belong to the Nikulin or Gritsenko--Cl\'ery automorphic catalogue
with different groups and multiplier systems; they lie outside the
primitive CHL rows for the $K3\times E$ host.

The Borcherds invariant is determined by $c_N(0)/2$, independently of
compact $\kch$ and fibre Euler terms.  At $N=1$ the Borcherds value is
$5$; the compact chiral supertrace and elliptic fibre Euler
characteristic give $0+0=0$.  The number $2$ is the K3 fibre witness
$\chi(\cO_{K3})$; the K3 fibre-rank invariant in the $K3\times E$
spectrum is the Mukai rank $24$.

\subsection*{Hall, BKM, and Mukai branches}

The stability-theoretic construction gives a chamberwise positive sector
for the critical Hall theory attached to $K3\times E$.  The BKM branch
is the primitive Borcherds denominator side, governed by $\Delta_5$.
The Hall-Drinfeld double enters after the Hall pairing and completion
are fixed.  The Mukai Yangian branch is a distinct K3 surface
construction governed by the Mukai lattice and self-mirror data.  Equal
graded characters or compatible scalar partition functions give
consistency evidence for comparison maps; algebra identifications
require bracket-level recognition data.

A Hall-to-Borcherds recognition theorem must identify the primitive Hall
generators with BKM roots, the Hall Euler form with the Borcherds
bilinear form, the reduced multiplicities with the Fourier coefficients
of the K3 elliptic genus, and the completed coproduct with the
Drinfeld-double coproduct.  Without these four identifications the
positive Hall sector, reduced DT scalar, and automorphic denominator
remain compatible but distinct objects.

\subsection*{Chamberwise construction}

The cited stability theorems provide the following conditional
construction.  Bridgeland stability on the K3 factor and Toda polynomial
stability on the product give a chamberwise selector for the positive
Hall sector of $K3\times E$, after orientation, compact-support, and
reduced data are fixed:
\[
  \Yp(K3\times E;\sigma)
  =
  \bigoplus_{\operatorname{Im}Z_\sigma(\gamma)>0}
  H^{*}\!\left(\Mum_\gamma^{\sigma\text{-ss}},\phi_W\right).
\]
Kontsevich-Soibelman wall crossing transports the completed ordered
Hall product between chambers.  Oberdieck-Pandharipande reduced DT
theory supplies the scalar trace
$-4096\,\Delta_5^{-2}$ in OP normalisation.  The separated values remain
\[
  \kcat(K3\times E)=0,\qquad
  \kch(K3\times E)=0,\qquad
  \kch^{\mathrm{Heis}}(K3\times E)=3,\qquad
  \kBKM(\Delta_5)=5,\qquad
  \kfib(K3\times E)=24.
\]
The constructions stay separate: stability chambers, reduced DT scalars,
Hall-Drinfeld doubles, and BKM denominators form distinct layers.  The
Mukai Yangian branch remains a separate K3 surface construction.
